% !TEX program = pdflatex
% Defensa TFM — Beamer (línea cronológica del documento)
\documentclass[aspectratio=169,10pt,professionalfonts]{beamer}

\usepackage[T1]{fontenc}
\usepackage[scaled=0.95]{helvet}
\renewcommand{\familydefault}{\sfdefault}
\usepackage{lmodern}
\renewcommand{\familydefault}{\sfdefault}
\usepackage[spanish]{babel}
\usepackage{amsmath,amssymb,bm}
\usepackage{booktabs}
\usepackage{tabularx}
\usepackage{colortbl}
\usepackage{graphicx}
\usepackage{tikz}
\usetikzlibrary{positioning,fit,arrows.meta,babel,calc}
\usepackage[authoryear,round]{natbib}

\graphicspath{{../Imagenes/}{../../POCs/data/output/figures/}{../}}

% ── Diseño alineado con la plantilla de defensa de UNIR (PowerPoint) ──
% Paleta de la plantilla: azul UNIR, azul oscuro, grises (sin verdes ni rojos).
\definecolor{unirazul}{HTML}{219BC3}
\definecolor{unirnavy}{HTML}{0E4A75}
\definecolor{unirgris}{HTML}{858591}
\definecolor{unirclaro}{HTML}{C0C0C8}
\definecolor{ink}{HTML}{3D3D4A}
% Los diagramas TikZ usan solo la paleta de la plantilla
\colorlet{blue}{unirazul}
\colorlet{orange}{unirnavy}
\colorlet{violet}{unirgris}

\usetheme{default}
\usefonttheme{professionalfonts}
\setbeamercolor{normal text}{fg=ink,bg=white}
\setbeamercolor{structure}{fg=unirazul}
\setbeamercolor{alerted text}{fg=unirnavy}
\setbeamercolor{block title}{bg=unirazul,fg=white}
\setbeamercolor{block body}{bg=unirazul!8}
\setbeamercolor{block title alerted}{bg=unirnavy,fg=white}
\setbeamercolor{block body alerted}{bg=unirnavy!8}
\setbeamerfont{block title}{series=\bfseries,size=\small}
\setbeamersize{text margin left=0.85cm,text margin right=0.85cm}
\setbeamertemplate{navigation symbols}{}
\setbeamertemplate{itemize item}{\textcolor{unirazul}{\small$\bullet$}}
\setbeamertemplate{itemize subitem}{\textcolor{unirazul}{\scriptsize$\circ$}}
\setbeamertemplate{enumerate items}[default]
\setbeamercolor{enumerate item}{fg=unirazul}
\addtobeamertemplate{itemize body begin}{}{\setlength{\itemsep}{1pt}\setlength{\parsep}{0pt}}
\addtobeamertemplate{enumerate body begin}{}{\setlength{\itemsep}{1pt}\setlength{\parsep}{0pt}}

% Título de diapositiva: línea azul (apartado) + línea verde azulada (título)
\makeatletter
\setbeamertemplate{frametitle}{%
  \vskip 0.30cm
  \edef\unir@sec{\detokenize\expandafter{\insertsectionhead}}%
  \edef\unir@ttl{\detokenize\expandafter{\insertframetitle}}%
  \ifx\unir@sec\unir@ttl
    {\fontsize{14}{16}\selectfont\bfseries\color{unirnavy}\insertframetitle\par}%
  \else
    {\fontsize{14}{16}\selectfont\bfseries\color{unirazul}\insertsectionhead\par}%
    {\fontsize{14}{16}\selectfont\bfseries\color{unirnavy}\insertframetitle\par}%
  \fi
  \vskip 0.12cm}
\makeatother
\setbeamerfont{frametitle}{size=\normalsize}

\tikzset{cuerpo/.style={rectangle, rounded corners=3pt, very thick, draw=#1, fill=#1!8,
                       align=left, inner sep=5pt, anchor=north}}

% Celdas de resultado de la síntesis Q vs C
\newcommand{\ganaQ}{\cellcolor{unirazul}\textcolor{white}{\textbf{Gana Q}}}
\newcommand{\ganaC}{\cellcolor{unirclaro}\textbf{Gana C}}
\newcommand{\empate}{\cellcolor{unirazul!14}Empate}

% Pie: logotipo UNIR a la izquierda y «pág. NN» con barra vertical a la derecha
\newcommand{\unirpagina}{%
  \begin{tikzpicture}[baseline=0pt]
    \draw[unirazul, line width=1.6pt] (0,-0.12) -- (0,0.50);
    \node[anchor=south west, font=\fontsize{8}{9}\selectfont\bfseries, text=unirclaro, align=left,
          inner sep=0pt] at (0.14,-0.12)
      {p\'ag.\\ \ifnum\insertframenumber<10 0\fi\insertframenumber};
  \end{tikzpicture}}
\setbeamertemplate{footline}{%
  \leavevmode\hbox to\paperwidth{%
    \hskip0.75cm\raisebox{0.10cm}{\includegraphics[height=0.62cm]{logo_unir.png}}%
    \hfill\raisebox{0.14cm}{\unirpagina}\hskip0.85cm}%
  \vskip0.12cm}

\title{Ventaja cuántica en ciencia de datos:\\
condiciones empíricas, métodos híbridos\\
y criterios de activación selectiva}
\subtitle{Defensa del Trabajo de Fin de Máster}
\author{Juan David Prada Malagón}
\institute{Máster en Computación Cuántica · UNIR\\[2pt]
Director: Ricardo Rafael Domínguez Reyes}
\date{}

\begin{document}

% =====================================================================
\section{Título}
% =====================================================================
\begin{frame}[plain]
\begin{tikzpicture}[remember picture,overlay]
  \fill[unirazul] (current page.north west) rectangle ($(current page.south east)+(0,1.15)$);
  \node[text=white, align=center, text width=11cm, inner sep=0pt] (tt) at ($(current page.center)+(0,0.55)$)
    {{\fontsize{20}{24}\selectfont\bfseries\hyphenpenalty=10000 Ventaja cuántica en ciencia de datos:\\
      condiciones empíricas,\\ métodos híbridos y criterios\\ de activación selectiva\par}
     \vspace{0.55cm}
     {\large Defensa del Trabajo de Fin de Máster\par}
     \vspace{0.3cm}
     {\normalsize Juan David Prada Malagón\par}
     \vspace{0.12cm}
     {\footnotesize Máster en Computación Cuántica $\cdot$ UNIR\\ Director: Ricardo Rafael Domínguez Reyes\par}};
  \draw[white, line width=1.1pt] ($(tt.north west)+(-0.55,0.45)$) -- ++(0,-0.7) ($(tt.north west)+(-0.55,0.45)$) -- ++(0.7,0);
  \draw[white, line width=1.1pt] ($(tt.south east)+(0.55,-0.45)$) -- ++(0,0.7) ($(tt.south east)+(0.55,-0.45)$) -- ++(-0.7,0);
  \node[anchor=south west, inner sep=0pt] at ($(current page.south west)+(0.75,0.12)$)
    {\includegraphics[height=0.82cm]{logo_unir.png}};
\end{tikzpicture}
\end{frame}

\begin{frame}[plain]
\begin{tikzpicture}[remember picture,overlay]
  \fill[unirazul] (current page.north west) rectangle (current page.south |- current page.south west);
  \fill[unirazul] (current page.north west) rectangle ($(current page.south)$);
  \node[text=white, anchor=west] at ($(current page.west)+(1.6,0)$) {{\fontsize{34}{38}\selectfont ÍNDICE}};
  \node[anchor=west, align=left, inner sep=0pt, font=\fontsize{11.5}{14}\bfseries\selectfont] at ($(current page.center)+(0.35,0)$)
    {
    \textcolor{unirclaro}{01 --}\,\textcolor{unirnavy}{Motivación}\\[0.30cm]
    \textcolor{unirclaro}{02 --}\,\textcolor{unirnavy}{Introducción}\\[0.30cm]
    \textcolor{unirclaro}{03 --}\,\textcolor{unirnavy}{Contexto y estado de la técnica}\\[0.30cm]
    \textcolor{unirclaro}{04 --}\,\textcolor{unirnavy}{Planteamiento y metodología}\\[0.30cm]
    \textcolor{unirclaro}{05 --}\,\textcolor{unirnavy}{Resultados: Q vs C}\\[0.30cm]
    \textcolor{unirclaro}{06 --}\,\textcolor{unirnavy}{Conjeturas: $\mathcal{Q}$ y QCR}\\[0.30cm]
    \textcolor{unirclaro}{07 --}\,\textcolor{unirnavy}{Evidencia de las conjeturas}\\[0.30cm]
    \textcolor{unirclaro}{08 --}\,\textcolor{unirnavy}{Conclusiones y limitaciones}\\[0.30cm]
    \textcolor{unirclaro}{09 --}\,\textcolor{unirnavy}{Trabajo futuro}
    };
\end{tikzpicture}
\end{frame}

% =====================================================================
\section{Motivación}
% =====================================================================
\begin{frame}{La brecha entre promesa y evidencia}
\begin{columns}[T]
\begin{column}{0.55\textwidth}
\begin{itemize}
  \item La ciencia de datos escala en dimensionalidad y complejidad combinatoria \citep{jordan2015, biamonte2017}.
  \item Los algoritmos variacionales y de kernel prometen ventaja en la era NISQ \citep{preskill2018, farhi2014, schuld2021}.
  \item Persiste una \textbf{brecha} entre promesas teóricas y evidencia reproducible sobre datos reales.
\end{itemize}
\end{column}
\begin{column}{0.42\textwidth}
\begin{block}{Pregunta práctica}
\textquestiondown{}Cuándo activar un módulo cuántico\\
y cuándo quedarse con el stack clásico?
\end{block}
\end{column}
\end{columns}
\end{frame}

\begin{frame}{Estructura del TFM}
\centering
\begin{tikzpicture}[font=\scriptsize, thick, scale=0.85, every node/.style={transform shape},
    capnodo/.style={rectangle, rounded corners=4pt, very thick,
                minimum width=2.8cm, minimum height=0.75cm,
                align=center, inner sep=3pt},
    bloque/.style={rectangle, rounded corners=6pt, dashed,
                   inner sep=6pt, very thick},
    arr/.style={->, thick, gray!60}]
  \node[capnodo, draw=blue!55!black, fill=blue!8]  (c1) at (0,    0) {\textbf{Cap. 1}\\Introducción};
  \node[capnodo, draw=blue!55!black, fill=blue!8]  (c2) at (0,  -1.35) {\textbf{Cap. 2}\\Contexto};
  \node[capnodo, draw=blue!55!black, fill=blue!8]  (c3) at (0,  -2.7) {\textbf{Cap. 3}\\Objetivos};
  \node[capnodo, draw=orange!60!black, fill=orange!7] (c4) at (4.6, -1.35) {\textbf{Cap. 4}\\Metodología};
  \node[capnodo, draw=orange!60!black, fill=orange!7] (c5) at (4.6, -2.7) {\textbf{Cap. 5}\\Evaluación};
  \node[capnodo, draw=violet!60!black, fill=violet!7] (c6) at (9.2, -0.8) {\textbf{Cap. 6}\\QCR / $\mathcal{Q}$};
  \node[capnodo, draw=violet!60!black, fill=violet!7] (c7) at (9.2, -2.15) {\textbf{Cap. 7}\\Conclusiones};
  \node[capnodo, draw=violet!60!black, fill=violet!7] (c8) at (9.2, -3.5) {\textbf{Cap. 8}\\Futuro};
  \node[bloque, draw=blue!40,   fit=(c1)(c2)(c3),
        label={[blue!60!black, font=\scriptsize\bfseries]above:Marco teórico}]   {};
  \node[bloque, draw=orange!45, fit=(c4)(c5),
        label={[orange!65!black, font=\scriptsize\bfseries]above:Implementación}] {};
  \node[bloque, draw=violet!45, fit=(c6)(c7)(c8),
        label={[violet!60!black, font=\scriptsize\bfseries]above:Discusión y cierre}] {};
  \draw[arr] (c3.east)  -- ++(0.55, 0) |- (c4.west);
  \draw[arr] (c4.south) -- (c5.north);
  \draw[arr] (c5.east)  -- ++(0.55, 0) |- (c6.west);
  \draw[arr] (c6.south) -- (c7.north);
  \draw[arr] (c7.south) -- (c8.north);
  \draw[->, thick, dashed, gray!50]
    (c5.north east) to[out=25, in=200] node[midway, above, font=\tiny, gray!65]{valida} (c6.south west);
\end{tikzpicture}
\end{frame}

% =====================================================================
\section{Introducción}
% =====================================================================
\begin{frame}{Planteamiento del problema}
\vspace{0.4em}
\begin{block}{Pregunta de investigación}
\itshape
\textquestiondown{}En qué condiciones, y para qué tipos de problemas de ciencia de datos,
los algoritmos cuánticos disponibles en la era NISQ ofrecen una
ventaja real y medible frente a los algoritmos clásicos?
\end{block}

\vspace{0.8em}
\begin{columns}[T]
\begin{column}{0.31\textwidth}
\centering\textbf{\textcolor{unirnavy}{Qué se compara}}\\[0.2em]
{\small algoritmos cuánticos frente al mejor clásico, en igualdad de condiciones}
\end{column}
\begin{column}{0.31\textwidth}
\centering\textbf{\textcolor{unirnavy}{Dónde}}\\[0.2em]
{\small cinco casos de uso comunes y tres clases de datos}
\end{column}
\begin{column}{0.31\textwidth}
\centering\textbf{\textcolor{unirnavy}{Qué se obtiene}}\\[0.2em]
{\small las condiciones en que gana el cuántico y un criterio para activarlo}
\end{column}
\end{columns}
\end{frame}

\begin{frame}{Alcance}
\begin{columns}[T]
\begin{column}{0.46\textwidth}
\textbf{Casos de uso comunes (5)}
\begin{itemize}
  \item Clasificación supervisada (QSVM)
  \item Regresión (QK-SVR)
  \item Detección de anomalías (OCSVM-ZZ)
  \item Optimización (QAOA / VQE)
  \item Series temporales (QK-SVR sobre lags)
\end{itemize}
{\footnotesize Nueve casos de uso con protocolo justo en el documento.}
\end{column}
\begin{column}{0.50\textwidth}
\textbf{Tres clases de datos}
\begin{enumerate}
  \item \textbf{Sintéticos con estructura cuántica nativa}: generados con el
        \emph{feature map} del kernel cuántico (ZZ); la etiqueta o el objetivo
        es el valor de un observable sobre el estado codificado
  \item \textbf{Práctica}: datos reales de SKG Tecnología
  \item \textbf{Referencia pública}
\end{enumerate}
\end{column}
\end{columns}
\end{frame}

% =====================================================================
\section{Contexto y estado de la técnica}
% =====================================================================
\begin{frame}{Era NISQ}
\begin{block}{NISQ \citep{preskill2018}}
Dispositivos de escala intermedia \textbf{ruidosos}: qubits limitados, decoherencia, sin corrección de errores a gran escala.
\end{block}
\begin{itemize}
  \item Algoritmos viables hoy: \textbf{variacionales} (QAOA, VQE) e híbridos \citep{farhi2014, cerezo2021}.
  \item Kernels cuánticos: acceso a espacios de Hilbert de alta dimensión \citep{schuld2021, huang2021}.
  \item Distinción clave: \textbf{metodologías} maduras vs.\ \textbf{hardware} aún inmaduro.
\end{itemize}
\end{frame}
% =====================================================================
\section{Planteamiento y metodología}
% =====================================================================
\begin{frame}{Metodología en cuatro fases}
\centering
\begin{tikzpicture}[font=\scriptsize,
    cab/.style={rectangle, rounded corners=3pt, minimum width=3.25cm, minimum height=0.7cm,
                align=center, text=white, font=\footnotesize\bfseries},
    arr/.style={->, very thick, >=stealth, unirnavy}]
  \node[cab, fill=unirazul] (h1) at (0,0) {1. Datos};
  \node[cab, fill=unirnavy] (h2) at (3.55,0) {2. Modelos};
  \node[cab, fill=unirgris] (h3) at (7.1,0) {3. Evaluación};
  \node[cab, fill=unirazul!70!unirnavy] (h4) at (10.65,0) {4. Híbrido y diagnóstico};
  \draw[arr] (h1) -- (h2); \draw[arr] (h2) -- (h3); \draw[arr] (h3) -- (h4);
  \node[cuerpo=unirazul, text width=2.85cm, minimum height=4.6cm] at (0,-0.5) {%
    \begin{itemize}\setlength{\itemsep}{3pt}\setlength{\leftmargini}{0.8em}
      \item Tres clases: sintéticos nativos (ZZ), práctica SKG, referencia pública
      \item Preprocesado: estandarizar, PCA a 2 componentes y escalar a $[-\pi,\pi]$
      \item Particiones aleatorias repetidas; en series, 70/30 en orden temporal
    \end{itemize}};
  \node[cuerpo=unirnavy, text width=2.85cm, minimum height=4.6cm] at (3.55,-0.5) {%
    \begin{itemize}\setlength{\itemsep}{3pt}\setlength{\leftmargini}{0.8em}
      \item \textbf{Cuánticos:} QSVM, QK-SVR y OCSVM con kernel ZZ de 2 qubits; QAOA y VQE
      \item \textbf{Clásicos:} SVM, regresión logística, Ridge, RF, GB, Isolation Forest, LOF;
            Hungarian, ILP, Greedy, NN+2-opt, SA-QUBO
    \end{itemize}};
  \node[cuerpo=unirgris, text width=2.85cm, minimum height=4.6cm] at (7.1,-0.5) {%
    \begin{itemize}\setlength{\itemsep}{3pt}\setlength{\leftmargini}{0.8em}
      \item \textbf{Protocolo justo:} tabla A (mismas 2 variables para todos) y tabla B (el clásico, con todas)
      \item \textbf{Métricas:} exactitud, $R^2$, AUC, $P$(óptimo)
      \item 5--10 repeticiones; simulación exacta del vector de estado
    \end{itemize}};
  \node[cuerpo=unirazul!70!unirnavy, text width=2.85cm, minimum height=4.6cm] at (10.65,-0.5) {%
    \begin{itemize}\setlength{\itemsep}{3pt}\setlength{\leftmargini}{0.8em}
      \item Diagnóstico $\mathcal{Q}(\mathcal{D})$: indicadores previos a ejecutar el algoritmo
      \item Criterios de activación del módulo cuántico
      \item QCR: mapeo hacia una representación cuántica (conjetura)
    \end{itemize}};
  \node at (5.3,-5.5) [rectangle, rounded corners=3pt, fill=unirnavy, text=white, font=\footnotesize\bfseries,
        minimum width=13.4cm, minimum height=0.7cm]
    {Salida: para cada caso de uso y clase de datos, \textquestiondown{}gana Q o gana C?};
\end{tikzpicture}
\end{frame}

\begin{frame}{Modelos comparados}
\centering
{\footnotesize
\renewcommand{\arraystretch}{1.2}
\begin{tabularx}{\linewidth}{@{}>{\raggedright\arraybackslash}p{1.9cm}>{\raggedright\arraybackslash}X>{\raggedright\arraybackslash}X@{}}
\toprule
\textbf{Caso de uso} & \textbf{Modelo cuántico} & \textbf{Modelos clásicos de referencia} \\
\midrule
Clasificación &
\textbf{QSVM-ZZ}: kernel de fidelidad con \texttt{ZZFeatureMap} (2 qubits, 2 repeticiones,
entrelazamiento completo) y SVM con matriz de Gram precalculada &
Gradient Boosting, SVM-RBF, SVM lineal, regresión logística \\
Regresión &
\textbf{QK-SVR-ZZ}: el mismo kernel con SVR ($C{=}10$, $\varepsilon{=}0.1$) &
Ridge, SVR-RBF, Random Forest, Gradient Boosting \\
Anomalías &
\textbf{OCSVM-ZZ}: SVM de una clase con el mismo kernel ($\nu{=}0.05$) &
OCSVM-RBF, Isolation Forest, LOF \\
Optimización &
\textbf{QAOA} (COBYLA; $p{=}1$ en asignación y TSP, $p{=}3$ en MDS y portafolio) y
\textbf{VQE} (\texttt{TwoLocal} Ry--CZ, 3 repeticiones) &
Hungarian / ILP (PuLP-CBC), Greedy, fuerza bruta, vecino más cercano + 2-opt, recocido simulado sobre el QUBO (SA-QUBO) \\
Series temporales &
\textbf{QK-SVR-ZZ} sobre los 2 últimos valores (lags) &
Ridge, SVR-RBF, Gradient Boosting \\
\bottomrule
\end{tabularx}}

\vspace{0.6em}
{\footnotesize Todos los modelos cuánticos se simulan con el vector de estado exacto (sin ruido de hardware).}
\end{frame}

\begin{frame}{Datasets utilizados}
\centering
{\footnotesize
\renewcommand{\arraystretch}{1.25}
\begin{tabularx}{\linewidth}{@{}>{\raggedright\arraybackslash}p{1.7cm}>{\raggedright\arraybackslash}X>{\raggedright\arraybackslash}X>{\raggedright\arraybackslash}X@{}}
\toprule
\textbf{Caso de uso} & \textbf{Sintético nativo} & \textbf{Práctica SKG} & \textbf{Referencia pública} \\
\midrule
Clasificación &
\textbf{ZZ sintético}: $x\in[-\pi,\pi]^2$; etiqueta = signo de $\langle Z\otimes Z\rangle$ &
\textbf{Incidentes NUSE (Bogotá)}: gravedad binaria del incidente &
\textbf{Telco Churn}: abandono de clientes (IBM / Kaggle) \\
Regresión &
\textbf{ZZ sintético}: objetivo = $\langle Z\otimes Z\rangle$ &
\textbf{Tiempos de viaje}: velocidad media a partir de lecturas y hora &
\textbf{Friedman \#1}: función no lineal sintética con ruido \\
Anomalías &
\textbf{\texttt{zz\_phase}}: la fase minoritaria del estado ZZ son los atípicos &
\textbf{Incidentes} con atípicos inyectados &
\textbf{\texttt{moons}}: dos lunas con atípicos uniformes en su caja \\
Series temporales &
\textbf{\texttt{zz\_lag}}: objetivo = $\langle Z\otimes Z\rangle$ de los 2 últimos valores &
\textbf{Velocidad de 5 corredores}: lecturas cada 5 minutos &
\textbf{\texttt{sine}}: senoide con dos estacionalidades y ruido \\
Optimización &
\textbf{Asignación} agente--tarea (QUBO) y \textbf{MDS} (red de 6 nodos y grafos Erdős--Rényi) &
Sin datos SKG (asignación real: trabajo futuro) &
\textbf{Portafolio} de 6 activos (Yahoo Finance) y \textbf{TSP} con 4 capitales de Colombia \\
\bottomrule
\end{tabularx}}
\end{frame}
% =====================================================================
\section{Resultados: Q vs C}
% =====================================================================
\begin{frame}{Mapa de errores: cuántico frente a clásico}
\centering
\includegraphics[width=\linewidth,height=0.68\textheight,keepaspectratio]{13_mapa_errores_q_vs_c.png}

\vspace{0.2em}
{\footnotesize El cuántico (azul) solo es el mejor con datos nativos; en optimización, QAOA iguala al exacto, no lo supera.}
\end{frame}

\begin{frame}{Síntesis: caso de uso $\times$ conjunto de datos}
\centering
{\scriptsize
\renewcommand{\arraystretch}{1.0}
\setlength{\tabcolsep}{4pt}
\begin{tabular}{@{}llll@{}}
\toprule
\textbf{Caso de uso} & \textbf{Clase de datos} & \textbf{Conjunto de datos} & \textbf{Resultado} \\
\midrule
Clasificación & Nativo & ZZ sintético & \ganaQ \\
 & SKG & Incidentes NUSE & \empate \\
 & Referencia & Telco Churn & \ganaC \\
\midrule
Regresión & Nativo & ZZ sintético & \ganaQ \\
 & SKG & Tiempos de viaje & \ganaC \\
 & Referencia & Friedman \#1 & \ganaC \\
\midrule
Anomalías & Nativo & \texttt{zz\_phase} & \ganaQ \\
 & SKG & Incidentes inyectados & \ganaC \\
 & Referencia & \texttt{moons} & \ganaC \\
\midrule
Series temporales & Nativo & \texttt{zz\_lag} & \ganaQ \\
 & SKG & Corredores & \ganaC \\
 & Referencia & \texttt{sine} & \ganaC \\
\midrule
Optimización & Nativo & Asignación QUBO & \empate \\
 & Nativo & MDS (grafos ER) & \ganaC \\
 & Referencia & Portafolio & \ganaC \\
 & Referencia & TSP $n{=}4$ & \empate \\
\bottomrule
\end{tabular}}

\vspace{0.1em}
{\scriptsize Empate: diferencia dentro de una desviación estándar de las repeticiones, o mismo valor óptimo.}
\end{frame}

\begin{frame}{Criterio de activación del módulo cuántico}
\centering
{\footnotesize
\renewcommand{\arraystretch}{1.15}
\begin{tabular}{@{}>{\raggedright\arraybackslash}p{2.6cm}>{\raggedright\arraybackslash}p{3.6cm}>{\raggedright\arraybackslash}p{3.2cm}>{\raggedright\arraybackslash}p{3.0cm}@{}}
\toprule
\textbf{Problema} & \textbf{Activar si} & \textbf{Módulo Q} & \textbf{Fallback} \\
\midrule
Clasif.\ /\ Regresión &
coincidencia mapa--generador
\emph{y} vence a GB/RF
(no basta $\mathcal{Q}>0$ vs RBF) &
QSVM / QK-SVR &
GB / RF \\
Anomalías &
novedad alineada al mapa
(\texttt{zz\_phase}) &
OCSVM-ZZ &
IF / LOF \\
Forecasting &
target alineado (\texttt{zz\_lag}) &
QK-SVR &
Ridge / GB \\
Optimización / TSP &
QUBO factible, $n$ viable &
QAOA / VQE &
Greedy / ILP / 2-opt \\
Caso general &
SKG/Telco sin estructura ZZ
o $n$ demasiado grande &
\textbf{no se activa} &
pipeline clásico \\
\bottomrule
\end{tabular}}
\end{frame}

% =====================================================================
\section{Conjeturas}
% =====================================================================
\begin{frame}{De la evidencia emergen dos preguntas}
\begin{block}{Preguntas originales (Cap.~6)}
\begin{enumerate}
  \item \textbf{P1 --- tests de cuanticidad}\\[0.35em]
        \textquestiondown{}Se puede diagnosticar \emph{a priori} si un dataset tiene
        \emph{estructura cuántica} (alineación mapa--dato) y, por tanto, admitirá ventaja?

  \vspace{1.0em}
  \item \textbf{P2 --- mapeos cuántico-compatibles}\\[0.35em]
        Si no hay estructura nativa, \textquestiondown{}existe un nuevo mapeo $\psi$ y un espacio
        de representación donde se conserve la estructura e información del dato,
        adquiera estructura cuántica, y haya mapa de vuelta si hace falta?
        (analogía con la transformada de Fourier)
\end{enumerate}
\end{block}
\end{frame}

\begin{frame}{Camino a la Conjetura 1 (I) --- por qué un perfil}
\begin{columns}[T]
\begin{column}{0.48\textwidth}
\begin{block}{Literatura}
\begin{itemize}
  \item Ventaja cuántica de kernel $\Rightarrow$
        alineación con etiquetas \textbf{y}
        geometría no reproducible por el
        clásico \citep{huang2021, liu2021}.
  \item Un solo indicador no basta:
        hace falta $A_Q > A_{cl}$
        \textbf{y} $d_F$ grande
        \citep{huang2021, schuld2021}.
  \item Riesgo: el kernel cuántico puede
        \textbf{concentrarse} y perder señal
        --- vigilar $H(K)$ \citep{thanasilp2022}.
\end{itemize}
\end{block}
\end{column}
\begin{column}{0.48\textwidth}
\begin{block}{Hecho empírico}
Ventaja si el mapa que estructura el dato
= mapa del modelo \textbf{cuántico}.
Si no: gana el \textbf{clásico}.
\end{block}
\vspace{0.5em}
\begin{block}{Idea corregida}
No un escalar: un \textbf{vector}
$\mathcal{Q}(\mathcal{D})=(\Delta A,\hat{d}_F,\Delta H,\hat g,\ldots)$
que se \emph{mira} componente a componente.
\end{block}
\end{column}
\end{columns}
\end{frame}

\begin{frame}{Camino a la Conjetura 1 (II) --- componentes y evidencia}
\begin{columns}[T]
\begin{column}{0.58\textwidth}
\begin{block}{Interpretación (cuántico vs clásico)}
{\footnotesize
\begin{itemize}
  \item $\Delta A$ --- \textquestiondown{}el cuántico ve mejor las clases?\\
        $>0$: sí, más alineación con las etiquetas
        que el clásico \citep{cortes2012, huang2021}.
  \item $\hat{d}_F$ --- \textquestiondown{}la geometría es distinta?\\
        Grande: el cuántico no es un clásico
        disfrazado \citep{liu2021}.
  \item $\Delta H$ --- \textquestiondown{}el kernel cuántico aporta
        diversidad o se ha ``aplanado''?\\
        Bajo/negativo: todas las similitudes se parecen
        (concentración) $\Rightarrow$ poca info útil
        \citep{thanasilp2022}.
  \item $\hat g$ --- \textquestiondown{}el RBF reproduce la geometr\'ia
        cu\'antica?\\
        Cerca de $1$: s\'i; grande: no (necesario, no suficiente)
        \citep{huang2021}.
\end{itemize}}
\end{block}
\end{column}
\begin{column}{0.40\textwidth}
\begin{block}{Evidencia (TFM)}
Barrido $\alpha$ ($n{=}9$, NB08):
mezcla de etiquetas, $X$ fijo.\\[0.25em]
Hallazgo robusto:
$\rho_S(\alpha,\Delta A){=}0.833$
($p{=}0.005$).\\[0.25em]
{\footnotesize $d_F$/$H$ constantes por diseño;
no confundir con
$\rho(\mathcal{Q},\mathrm{ventaja})$.}
\end{block}
\end{column}
\end{columns}
\end{frame}

\begin{frame}{Conjetura 1 --- $\mathcal{Q}(\mathcal{D})$ como \textbf{vector}, no escalar}
\begin{equation}
  \mathcal{Q}(\mathcal{D})
  \;=\;
  \bigl(\Delta A,\; \hat{d}_F,\; \Delta H,\; \hat g\bigr)
  \qquad\text{(perfil; sin pesos)}
  \label{eq:Q}
\end{equation}
{\footnotesize
Opcional en el documento: un escalar $w^\top\mathcal{Q}$ se probó y \textbf{falla} como ranking
operativo --- por eso no se usa para decidir.}

\vspace{0.4em}
\begin{columns}[T]
\begin{column}{0.48\textwidth}
\begin{block}{Cómo se lee}
{\footnotesize
\begin{itemize}
  \item $\Delta A$: \textquestiondown{}mejor alineación
        que el clásico?
  \item $\hat{d}_F$: \textquestiondown{}geometría distinta?
  \item $\Delta H$: \textquestiondown{}diversidad o
        concentración?
  \item $\hat g$: \textquestiondown{}el RBF reproduce
        la geometr\'ia cu\'antica?
\end{itemize}}
\end{block}
\end{column}
\begin{column}{0.48\textwidth}
\begin{alertblock}{Límite}
{\footnotesize
El vector describe geometría vs.\ RBF.
La adopción exige ventaja vs.\
mejor clásico \textbf{en el caso de uso}
(clasif./reg./anomalías) y alineación
mapa--generador.}
\end{alertblock}
\end{column}
\end{columns}
\end{frame}

\begin{frame}{Validación P1 --- barrido sintético ($\alpha$)}
\begin{columns}[T]
\begin{column}{0.62\textwidth}
\centering
\includegraphics[width=\linewidth,height=0.68\textheight,keepaspectratio]{08_barrido_sintetico_alpha.png}
\end{column}
\begin{column}{0.36\textwidth}
\begin{block}{Qué es (NB08 real)}
{\footnotesize
\begin{itemize}
  \item $\alpha$: fracción de etiquetas
        cuánticas ($X$ fijo).
  \item $\Delta A$ crece con $\alpha$:
        $\rho_S{=}0.833$ ($p{=}0.005$).
  \item $d_F$ y $H$ \textbf{constantes}
        (kernels no cambian con labels).
\end{itemize}}
\end{block}
\begin{alertblock}{Límite}
{\footnotesize
$\rho{=}0.833$ es $\rho(\alpha,\Delta A)$,
no $\rho(\mathcal{Q},\text{ventaja})$.
Barrido structure-aware: trabajo futuro.}
\end{alertblock}
\end{column}
\end{columns}
\end{frame}

\begin{frame}{P1 --- qué comparamos (mismas métricas)}
\begin{block}{Estructura cuántica = vector $(\Delta A,\ \Delta H,\ d_F,\ g,\ \mathcal{Q})$}
{\footnotesize
\begin{itemize}
  \item $\Delta A = A(K_Q,yy^\top)-A(K_{\mathrm{RBF}},yy^\top)$ --- alineaci\'on de \citep{cortes2012}
  \item $\Delta H = H(K_Q)-H(K_{\mathrm{RBF}})$ --- proxy de \citep{thanasilp2022}
  \item $d_F=\|K_Q-K_{\mathrm{RBF}}\|_F/n$ --- proxy propio (cf.\ \citep{liu2021})
  \item $g$ (diferencia geom.) --- \citep{huang2021}; $\mathcal{Q}$ compuesto (este TFM)
\end{itemize}
\textbf{Mismo protocolo} en clasificación, regresión, anomalías y optimización
(ZZ 2q, $n{=}60$, 3 repeats, RBF Frobenius-best).\\
Optimización = \textbf{paisaje QUBO} (bitstrings $\to$ energía), no QAOA.\\
Control nulo: $\Delta A$ con $y$ permutada. Sin accuracy / $R^2$ / AUC / coste.}
\end{block}
\vspace{0.2em}
\begin{alertblock}{Nota: no hay $\Delta_{\mathrm{op}}$ aquí}
{\footnotesize
Eso era solo un nombre confuso para
``accuracy(QSVM)$-$accuracy(mejor clásico)'' en clasificación.
\textbf{No} es métrica de estructura. Se retiró de estas figuras.}
\end{alertblock}
\end{frame}

\begin{frame}{P1 --- qué mide $g$ (Huang et al., 2021)}
\begin{columns}[T]
\begin{column}{0.56\textwidth}
\centering
\includegraphics[width=\linewidth,height=0.56\textheight,keepaspectratio]{huang2021_fig2_geometria_kernel.png}\\[-0.2em]
{\tiny Fig.~2 de \citet{huang2021}, CC~BY~4.0}
\end{column}
\begin{column}{0.42\textwidth}
\begin{block}{Idea}
{\footnotesize
$g$ compara las \textbf{flechas} (similitudes kernel) entre puntos en el espacio
cl\'asico y en el cu\'antico. $g$ peque\~no: el kernel cl\'asico ya reproduce
la geometr\'ia cu\'antica; $g$ grande: no.\\[0.3em]
$d$ = dimensi\'on efectiva del dato en Hilbert.}
\end{block}
\begin{alertblock}{Nuestra medici\'on}
{\scriptsize
$g$: necesario, no suficiente; no usa $y$.
Huang ($\lambda{=}10^{-3}$): $g\approx 76$--$85$ en casi todos los casos
(42 en anomal\'ias-SKG, artefacto). $g\gg1$ pero igual para nativos y reales:
\textbf{no discrimina}. Antes $\approx 3000$ (artefacto).}
\end{alertblock}
\end{column}
\end{columns}
\end{frame}

\begin{frame}{P1 --- los cinco indicadores en los cuatro casos de uso}
\begin{columns}[T]
\begin{column}{0.66\textwidth}
\centering
\includegraphics[width=\linewidth,height=0.70\textheight,keepaspectratio]{08_estructura_comparativa_casos_uso.png}
\end{column}
\begin{column}{0.32\textwidth}
\begin{block}{Cómo leerlo}
{\footnotesize
Un panel por indicador; eje X = caso de uso;
3 barras = Nativo / SKG / Referencia
(optimizaci\'on: asignaci\'on / MaxCut / TSP).\\[0.3em]
\begin{itemize}
  \item $\Delta A$: solo los nativos superan el ruido (banda azul);
        el resto, dentro del ruido o por debajo (incluida la optimizaci\'on)
  \item $\Delta H$, $d_F$, $g$: casi planos;
        excepci\'on anomal\'ias-SKG = artefacto de escala (outliers $8\sigma$ comprimen el MinMax)
  \item $\mathcal{Q}$: informativo; lo dominan indicadores planos
\end{itemize}}
\end{block}
\end{column}
\end{columns}
\end{frame}

\begin{frame}{P1 --- cómo interpretar}
\begin{block}{\textquestiondown{}$\Delta A$ negativo + $\Delta H$ mínimo = estructura?}
{\footnotesize
\textbf{No.} Los \textbf{nativos} tienen $\Delta A>0$ sobre el ruido
(clasif.\ $+0.20$, reg.\ $+0.28$, anom.\ $+0.09$); los datos reales y los
paisajes de optimizaci\'on quedan dentro o por debajo del nulo
($\Delta A<0$ = el RBF alinea mejor, no ``estructura negativa'').
El \'unico $\Delta H$ m\'inimo (anomal\'ias-SKG, $0.44$) es un artefacto de escala.}
\end{block}
\vspace{0.35em}
\begin{block}{\textquestiondown{}$d_F$ discrimina?}
{\footnotesize
\textbf{Aqu\'i, apenas}: $\approx 0.01$ en todos los casos; $g$ y $\Delta H$ tambi\'en planos.
Solo $\Delta A$ usa $y$; los otros tres solo ven $X$.}
\end{block}
\vspace{0.35em}
\begin{alertblock}{Lectura defendible}
{\footnotesize
Lo que separa estructura es $\Delta A$ (alineaci\'on mapa--dato \citep{liu2021,huang2021}),
no $d_F$, $g$ ni $\Delta H$. Optimizaci\'on: paisaje QUBO con 2 qubits
(PCA retiene $\approx 1/3$): \textbf{sin evidencia} de estructura; l\'imite del dise\~no.}
\end{alertblock}
\end{frame}

\begin{frame}{Camino a la Conjetura 2 --- analogía con Fourier}
\begin{block}{La idea, por analogía con la transformada de Fourier}
Un problema difícil en un espacio puede volverse más simple en otro;
allí se resuelve y, si hace falta, se regresa. Eso funciona gracias a tres propiedades:
\begin{enumerate}
  \item \textbf{conserva la información} (es invertible),
  \item \textbf{hace visible la estructura} del problema,
  \item tiene un \textbf{mapa de vuelta}.
\end{enumerate}
\end{block}
\vspace{0.2em}
\begin{alertblock}{La analogía con el kernel cuántico}
Si el dato no tiene estructura cuántica, \textquestiondown{}existe un \textbf{mapeo $\psi$ y un espacio de
representación $\mathcal{Z}$} donde se conserve la información del dato, adquiera estructura cuántica y se
pueda sacar partido de los algoritmos cuánticos, con un mapeo de vuelta al espacio original cuando haga falta?
\end{alertblock}
\end{frame}

\begin{frame}{Conjetura 2 --- espacio de representación QCR}
\begin{center}
\begin{tikzpicture}[font=\footnotesize, thick,
    box/.style={rectangle, rounded corners=3pt, draw=blue!55!black, fill=blue!8,
                minimum width=2.4cm, minimum height=0.7cm, align=center}]
  \node[box] (X) at (0,0) {Espacio original\\$\mathcal{X}$};
  \node[box, minimum width=3.6cm] (Z) at (5.2,0) {Representación cuántica\\$\mathcal{Z}$};
  \node[box] (Y) at (10.2,0) {Resultado\\$\hat y$ o solución};
  \draw[->, very thick] (X) -- node[above]{$\psi$} (Z);
  \draw[->, very thick] (Z) -- node[above]{$f_Q$} (Y);
  \draw[->, dashed, thick] (Y.south) to[bend left=10] node[below]{$\rho=\psi^{-1}$ (si hace falta)} (X.south);
\end{tikzpicture}
\end{center}
\vspace{-2.2em}
\begin{columns}[T]
\begin{column}{0.50\textwidth}
\begin{block}{Qué se pide de $\psi$}
{\footnotesize
\begin{enumerate}
  \item $\psi:\mathcal{X}\to\mathcal{Z}$ es un \textbf{isomorfismo}: conserva estructura e información.
  \item La \textbf{imagen} de $\psi$ tiene \textbf{estructura cuántica}.
  \item El mapa de vuelta es su inversa, $\rho=\psi^{-1}$.
\end{enumerate}
}
\end{block}
\end{column}
\begin{column}{0.47\textwidth}
\begin{alertblock}{Comentarios}
{\footnotesize
\begin{itemize}
  \item No sabemos con certeza la forma de $\psi$.
  \item El coste de calcular $\psi$ debe tenerse en cuenta en las aplicaciones.
  \item Queda como \textbf{conjetura abierta}.
\end{itemize}
}
\end{alertblock}
\end{column}
\end{columns}
\end{frame}

% =====================================================================
\section{Conclusiones}
% =====================================================================
\begin{frame}{Mensajes clave}
\begin{columns}[T]
\begin{column}{0.48\textwidth}
\begin{block}{1. Ventaja condicionada}
Solo gana el cuántico cuando el \emph{feature map} coincide con el generador de los datos.
En datos reales (SKG, Telco) gana el clásico.
\end{block}
\begin{block}{2. Optimización: valor de modelado}
QAOA iguala al óptimo exacto en problemas pequeños, pero no lo supera.
\end{block}
\end{column}
\begin{column}{0.48\textwidth}
\begin{block}{3. Diagnosticar antes de invertir}
$\mathcal{Q}(\mathcal{D})$, como vector de indicadores, orienta la decisión; QCR queda planteada como conjetura.
\end{block}
\begin{block}{4. Adopción selectiva e híbrida}
El módulo cuántico complementa al stack clásico; no lo sustituye.
\end{block}
\end{column}
\end{columns}
\vspace{0.3em}
\begin{alertblock}{Respuesta a la pregunta del TFM}
Hay ventaja medible cuando existe coincidencia mapa--generador; fuera de ahí, el clásico domina, y eso también es un resultado.
\end{alertblock}
\end{frame}

\section{Limitaciones}

\begin{frame}{Limitaciones del estudio}
\begin{columns}[T]
\begin{column}{0.48\textwidth}
\begin{block}{Simulación ideal}
Statevector sin ruido de hardware; el coste del kernel crece como $O(n^2)$.
\end{block}
\begin{block}{Escala reducida}
2 qubits tras PCA y muestras pequeñas: los resultados son orientativos.
\end{block}
\end{column}
\begin{column}{0.48\textwidth}
\begin{block}{Alcance de los datos}
SKG se usó como práctica (series: 5 corredores, un paso, sin ARIMA ni Prophet).
\end{block}
\begin{block}{Base teórica abierta}
$\mathcal{Q}(\mathcal{D})$ es empírico y QCR carece de demostración.
\end{block}
\end{column}
\end{columns}
\end{frame}

\section{Trabajo futuro}

\begin{frame}{Hoja de ruta}
\centering
\begin{tikzpicture}[font=\footnotesize,
    fase/.style={rectangle, rounded corners=3pt, minimum width=4.3cm, minimum height=0.8cm,
                 align=center, text=white, font=\footnotesize\bfseries},
    arr/.style={->, line width=1.6pt, >=stealth, unirnavy}]
  \node[fase, fill=unirazul] (h1) at (0,0) {Corto plazo};
  \node[fase, fill=unirnavy] (h2) at (4.9,0) {Medio plazo};
  \node[fase, fill=unirgris] (h3) at (9.8,0) {Largo plazo};
  \draw[arr] (h1) -- (h2); \draw[arr] (h2) -- (h3);
  \node[cuerpo=unirazul, text width=4.0cm, minimum height=3.4cm] at (0,-0.6) {%
    \begin{itemize}\setlength{\itemsep}{3pt}\setlength{\leftmargini}{0.8em}
      \item Hardware real con mitigación de errores
      \item Más qubits y menos reducción por PCA
      \item Más \emph{feature maps}
    \end{itemize}};
  \node[cuerpo=unirnavy, text width=4.0cm, minimum height=3.4cm] at (4.9,-0.6) {%
    \begin{itemize}\setlength{\itemsep}{3pt}\setlength{\leftmargini}{0.8em}
      \item Validar $\mathcal{Q}(\mathcal{D})$ en un corpus amplio
      \item QAOA sobre asignación y rutas reales de SKG
      \item Pruebas A/B cuántico--clásico
    \end{itemize}};
  \node[cuerpo=unirgris, text width=4.0cm, minimum height=3.4cm] at (9.8,-0.6) {%
    \begin{itemize}\setlength{\itemsep}{3pt}\setlength{\leftmargini}{0.8em}
      \item Demostrar o refutar QCR
      \item Benchmarks con separación formal
      \item Hardware tolerante a fallos
    \end{itemize}};
\end{tikzpicture}
\end{frame}

\begin{frame}[plain]
\begin{tikzpicture}[remember picture,overlay]
  \node[align=center, inner sep=0pt, text width=8cm] (g) at (current page.center)
    {{\fontsize{30}{34}\selectfont\bfseries\textcolor{unirazul}{muchas} \textcolor{unirnavy}{gracias}}\par
     \vspace{0.35cm}{\large\textcolor{unirnavy}{\textquestiondown{}Preguntas?}}\par
     \vspace{0.3cm}{\small\textcolor{unirgris}{Juan David Prada Malagón $\cdot$ UNIR}}};
  \draw[unirazul, line width=1.3pt] ($(g.north west)+(-1.0,0.6)$) -- ++(0,-0.9) ($(g.north west)+(-1.0,0.6)$) -- ++(0.9,0);
  \draw[unirazul, line width=1.3pt] ($(g.south east)+(1.0,-0.6)$) -- ++(0,0.9) ($(g.south east)+(1.0,-0.6)$) -- ++(-0.9,0);
  \node[anchor=south west, inner sep=0pt] at ($(current page.south west)+(0.75,0.12)$)
    {\includegraphics[height=0.82cm]{logo_unir.png}};
\end{tikzpicture}
\end{frame}

\section{Referencias}

\begin{frame}[allowframebreaks]{Selección}
\bibliographystyle{plainnat}
\begin{thebibliography}{99}
\bibitem[Biamonte et~al.(2017)]{biamonte2017}
Biamonte, J.\ et~al. (2017). \emph{Quantum machine learning}. Nature, 549, 195--202.

\bibitem[Cortes et~al.(2012)]{cortes2012}
Cortes, C., Mohri, M., \& Rostamizadeh, A. (2012). \emph{Algorithms for learning kernels based on centered alignment}. JMLR, 13, 795--828.

\bibitem[Cerezo et~al.(2021)]{cerezo2021}
Cerezo, M.\ et~al. (2021). \emph{Variational quantum algorithms}. Nat.\ Rev.\ Phys., 3, 625--644.

\bibitem[Farhi et~al.(2014)]{farhi2014}
Farhi, E., Goldstone, J., \& Gutmann, S. (2014). \emph{A quantum approximate optimization algorithm}. arXiv:1411.4028.

\bibitem[Fastovets et~al.(2019)]{fastovets2019}
Fastovets, D.\ V.\ et~al. (2019). \emph{Machine learning methods in quantum computing theory}. arXiv:1906.10175.

\bibitem[Huang et~al.(2021)]{huang2021}
Huang, H.-Y.\ et~al. (2021). \emph{Power of data in quantum machine learning}. Nat.\ Commun., 12, 2631.

\bibitem[Jordan \& Mitchell(2015)]{jordan2015}
Jordan, M.\ I., \& Mitchell, T.\ M. (2015). \emph{Machine learning: Trends, perspectives, and prospects}. Science, 349(6245), 255--260.

\bibitem[Liu et~al.(2021)]{liu2021}
Liu, Y., Arunachalam, S., \& Temme, K. (2021). \emph{A rigorous and robust quantum speed-up in supervised machine learning}. Nat.\ Phys., 17, 1013--1017.

\bibitem[Preskill(2018)]{preskill2018}
Preskill, J. (2018). \emph{Quantum computing in the NISQ era and beyond}. Quantum, 2, 79.

\bibitem[Schuld(2021)]{schuld2021}
Schuld, M. (2021). \emph{Supervised quantum machine learning models are kernel methods}. arXiv:2101.11020.

\bibitem[Thanasilp et~al.(2024)]{thanasilp2022}
Thanasilp, S., Wang, S., Cerezo, M., \& Holmes, Z. (2024). \emph{Exponential concentration in quantum kernel methods}. Nat.\ Commun., 15, 5200.
\end{thebibliography}
\end{frame}

\end{document}
